\documentclass[12pt, oneside, openany]{book}

\usepackage{mathptmx} % Contiene una fuente similar a Times New Roman

\usepackage[spanish, es-tabla]{babel} % Permite escritura en castellano
\usepackage[utf8]{inputenc} % Permite utilizar caracteres UTF8

\usepackage{graphicx} % Para la inclusión de gráficos e imágenes
\graphicspath{ {images/} } % Ruta para buscar las imágenes
\usepackage[a4paper,top=30mm,left=30mm,right=25mm,bottom=25mm,headheight=20mm]{geometry} % Configuración de los margenes de la página

% Paquetes para que funcione el formato.
\usepackage{titlesec}
\usepackage{setspace}
\usepackage{ragged2e}
\usepackage{fancyhdr}
\usepackage{lastpage}
\usepackage{stackengine}
\usepackage{array}
\usepackage[hidelinks]{hyperref}
\usepackage{enumitem}
\usepackage{float}
\usepackage{hypcap}
\usepackage{caption}
\usepackage{fancyvrb}
\usepackage{hyperref} % Paquete para que las referencias funcionen, y permite introducir links
\usepackage{xcolor} % Paquete para trabajar con colores (fondo de celdas, color del texto...)

% Se define un color gris desde su código RGB
\definecolor{gris}{RGB}{220,220,220}

\setcounter{secnumdepth}{3} % Para permitir numerar las sub-subsecciones

% Modifica el nombre de los índices al castellano
\addto\captionsspanish{
  \renewcommand{\contentsname}{Índice de contenido}
  \renewcommand{\listfigurename}{Índice de figuras}
  \renewcommand{\listtablename}{Índice de tablas}
}

% Formateo de los nombres de los apartados:
\titleformat{\chapter}[block]
  {\normalfont\Huge\bfseries\singlespacing}{\thechapter.}{1em}{\Huge}
\titlespacing*{\chapter}{0pt}{-62pt}{0pt}

\titleformat{\section}[block]
  {\normalfont\Large\bfseries}{\thesection.}{4pt}{\Large}
\titlespacing*{\section}{0pt}{\baselineskip}{0pt}

\titleformat{\subsection}[block]
  {\normalfont\large\bfseries}{\thesubsection.}{4pt}{\normalsize\large}
\titlespacing*{\subsection}{0pt}{0pt}{0pt}

\titleformat{\subsubsection}[block]
  {\normalfont\normalsize\bfseries}{\thesubsubsection.}{4pt}{\normalsize}
\titlespacing*{\subsubsection}{0pt}{0pt}{0pt}

\def\tablename{Tabla}

%% Variables para portada y cabeceras
%% Cambiar los valores para cada documento!!!
\def\title{Primera entrega de prácticas}
\def\subject{Pruebas y Despliegue de Software}
\def\author{Juan Francisco Mier Montoto}
\def\target{www.mier.info}
\def\authorid{UO283319}
\def\date{enero 2024}
\def\org{Escuela Politécnica de Ingeniería de Gijón}
\def\area{Grado en Ingeniería Informática en Tecnologías de la Información}

\def\ORG{\expandafter\MakeUppercase\expandafter{\org}}
\def\AREA{\expandafter\MakeUppercase\expandafter{\area}}
\def\SUBJECT{\expandafter\MakeUppercase\expandafter{\subject}}

\captionsetup{justification=centering}
\setlength{\headheight}{65pt}

\fancyhf{}
\fancyhead[L]{\includegraphics[height=16mm]{style/square.png}
  \hspace{1em} \Longstack[l] {
    \textbf{\SUBJECT} \newline
    \textbf{\title}}
  \newline \leftmark{}
}
\fancyhead[R]{\bfseries{Hoja \hyperlink{toc}{\thepage}~de~\pageref{LastPage}}}
\fancyfoot[C]{\href{\target}{\author}}
\renewcommand{\headrulewidth}{0pt} % default is 0pt
\renewcommand{\footrulewidth}{0.4pt} % default is 0

\fancypagestyle{plain}{%
  \fancyhf{}
  \fancyhead[L]{\includegraphics[height=16mm]{style/square.png}
    \hspace{1em} \Longstack[l]{
      \textbf{\SUBJECT} \newline
      \textbf{\title}}}
  \fancyhead[R]{\bfseries{Hoja \hyperlink{toc}{\thepage}~de~\pageref{LastPage}}}
  \fancyfoot[C]{\href{\target}{\author}}
  \renewcommand{\headrulewidth}{0pt} % default is 0pt
  \renewcommand{\footrulewidth}{0.4pt} % default is 0pt
}

\pagestyle{fancy}

\restylefloat{table}


\begin{document}

\rmfamily % Fuente tipo Romana

% Portada de la memoria
\begin{titlepage}
    \centering
    \bfseries {
        \null{}
        \vspace{0cm}
        \begin{table}[h]
            \centering
            \begin{tabular}{m{10cm} m{1cm} m{3cm}}
                \vspace{0.2cm}
                \includegraphics[width=86mm]{style/full.png} &  & \vspace{1.52mm} \includegraphics[width=23mm]{style/square.png} \\
            \end{tabular}
        \end{table}

        \vspace{3\baselineskip}

        \Large{\ORG{} \\ \vspace{3\baselineskip}}
        \large {
            \AREA{} \\ \vspace{3\baselineskip}
            \subject{} \\ \vspace{2\baselineskip}

            % TRABAJO FIN DE GRADO/MÁSTER Nº XXXXXXXXX \vspace{\baselineskip} \\
            \title{} \\ \vspace{3\baselineskip}

            \author{} \\
            \authorid{} \\
            % TUTOR/ES: \\
            % D. APELLIDO1 APELLIDO2, Nombre \\
            % D. APELLIDO1 APELLIDO2, Nombre \\  \vspace{\baselineskip}

            \vspace{2\baselineskip}
            FECHA:\@ \date{}
        }
    }
\end{titlepage}


% Índice de contenido
\addcontentsline{toc}{chapter}{Índice de contenido} % Añade la referencia al índice de contenido
\hypertarget{toc}{}
\tableofcontents

% Índice de figuras
\addcontentsline{toc}{chapter}{Índice de figuras}  % Añade la referencia al índice de contenido
\hypertarget{lof}{}
\listoffigures

\justify{} % Texto justificado
\setlength{\parskip}{\baselineskip} % Separación entre párrafos de 1 linea
\onehalfspacing{}

%% El contenido de la memoria, dividido en capítulos:
\chapter{Introducción}
Durante el desarrollo de software, es fundamental asegurarse de que el producto
final cumple con los requisitos y expectativas del cliente. Para ello, es
necesario realizar pruebas exhaustivas que permitan detectar fallos y errores
antes de que el producto llegue a manos del cliente.

En este documento se detalla el proceso de diseño de pruebas realizado para una
aplicación de juguete, que permite calcular el tipo de triángulo que se forma
según la longitud de los lados introducidos. Se detallan los casos de prueba
diseñados, así como los fallos detectados y su severidad.

\input{sections/2_diseño.tex}
\chapter{Reporte de fallos}
% ordenados por severidad
Tras la ejecución del diseño de pruebas anterior, se obtiene la siguiente lista
de fallos, ordenados por severidad. La \textit{severidad} que se le asigna a
cada fallo es totalmente subjetiva y se basa en la importancia que tiene el
mismo en el correcto funcionamiento de la aplicación.

Primero de todo, se ha de tener en cuenta los fallos de truncamiento que se producen
al introducir números decimales en los campos de entrada, algo fundamental que
provoca salidas erróneas.

También se considera grave el fallo que causa el incumplimiento de la desigualdad
triangular cuando la suma de los dos últimos lados es inferior que el primero.
Este fallo lógico se considera el grave ya que incumple claramente una de las
propiedades fundamentales de los triángulos.

El fallo que provoca el cierre inesperado de la aplicación al introducir un
caracter alfabético en uno de los campos debería corregirse con cierta urgencia,
puesto que es un detalle esencial que se debería haber comprobado durante el
desarrollo y que cualquier usuario puede cometer fácilmente.

Otro de los detalles que se deberían haber tenido en cuenta es el resultado de
introducir tres ceros en los campos de entrada, ya que se produce ``equilátero''
cuando, evidentemente, se trata de un triángulo inválido.

El hecho de que la aplicación considere como ``rectángulos'' a los triágulos que
los son no es un fallo lógico ni técnico, pero era un detalle que se debería haber
estipulado durante la fase de diseño.

\chapter{Resumen ejecutivo}
Durante la ejecución de los casos de prueba, se han detectado los cuatro fallos
estipulados en el reporte de fallos anterior, y se comprueba que la aplicación
cuenta con casos no tan extremos que no se habían tenido en cuenta que afectan
al correcto funcionamiento de la misma.

Puesto que se desconoce la implementación del código, se ha de asumir por la
complejidad de la aplicación y las características de los fallos que no son
defectos difíciles de corregir, pero sí que requieren un cambio en la lógica
de desarrollo teniendo en cuenta la importancia de dichos casos límite.

La detección de estos fallos ha sido posible gracias al correcto diseño de
pruebas realizado en clase.


\end{document}
